\section*{A tensor-prime quadratic obstruction to every fixed good-coordinate GU descent}

Let $n\ge2$ be even and define the symmetric Boolean phase
\[
 q_n(x)=(-1)^{\sum_{1\le r<s\le n}x_rx_s}
       =(-1)^{\binom{|x|}{2}}.                         \tag{1}
\]
This family is genuine Boolean and uses every physical coordinate.

\subsection*{Exact spectrum and GU values}

The polar matrix of the quadratic form in (1) is the matrix with zero
diagonal and ones off the diagonal.  Over $\mathbb F_2$ it has full rank
when $n$ is even.  Hence $q_n$ is bent and
\[
 \widehat q_n(S)^2=2^{-n}\quad(S\subseteq[n]).          \tag{2}
\]
For a fixed $S$, (2) gives $\pi(S)=2^{-|S|}$.  Therefore
\[
 U(q_n)=\mathbb E_{S\sim\mathrm{Unif}(2^{[n]})}(n-|S|)
       ={n\over2},\qquad I(q_n)={n\over2}.             \tag{3}
\]

To compute physical children, put $m=n-1$, which is odd.  The two
restrictions are $q_m$ and $q_m\chi_{[m]}$.  The polar form in odd dimension
has kernel spanned by the all-ones vector.  Consequently their squared
spectra are uniform, with mass $2^{-(m-1)}$, on the two opposite parity
hyperplanes of subsets.

For such a parity-hyperplane law and a supported set $S$ of size $k<m$,
exactly $2^{m-k-1}$ supported supersets contain $S$, so $\pi(S)=2^{-k}$.
For the hyperplane containing $[m]$, the top set instead has
$\pi([m])=2^{-(m-1)}$.  It follows that the two child $U$ values are
\[
 {m-2\over2},\qquad {m-2\over2}+2^{-(m-1)}.            \tag{4}
\]
Their own one-coordinate children are even-dimensional bent phases, all
with $U=(m-1)/2$.  Thus their $G$ values are
\[
 -{m\over2},\qquad -{m\over2}+m2^{-(m-1)}.             \tag{5}
\]
Equations (3)--(5) give
\[
 G(q_n)={3n\over2}-n2^{1-n},                           \tag{6}
\]
and, for every coordinate $i$,
\[
 \operatorname{Inf}_i(q_n)={1\over2},\qquad
 \Delta_iU={3\over2}-2^{1-n},                         \tag{7}
\]
\[
 \Delta_iG={4n-1\over2}-(2n-1)2^{1-n}.                \tag{8}
\]

In particular
\[
 {\Delta_iG\over\operatorname{Inf}_i}
 =4n-1-(2n-1)2^{2-n}\longrightarrow\infty.            \tag{9}
\]
At $n=4$, (9) is $53/4$, recovering the earlier finite gate.  At $n=6$ it
is $357/16=22.3125$, already refuting the proposed constant $14$.

More generally, the left side of the good-coordinate criterion for a fixed
pair $(a,b)$, divided by influence, is
\[
 {\Delta_iG-b\Delta_iU\over\operatorname{Inf}_i}
 =4n-1-(2n-1)2^{2-n}-b(3-2^{2-n}).                    \tag{10}
\]
For every fixed finite $a$ and $0\le b<1$, (10) exceeds $a$ for all
sufficiently large even $n$.  Thus no universal GU pair can be proved by
requiring the one-coordinate residual reversal at every tensor-prime
function.

\subsection*{Tensor primeness and scope}

The family is tensor-prime.  If a nontrivial disjoint factorization separated
coordinates $i$ and $j$, the multiplicative rectangle product at the zero
input would be $+1$.  For (1) it is
\[
 q_n(0)q_n(e_i)q_n(e_j)q_n(e_i+e_j)=-1,               \tag{11}
\]
because the mixed quadratic derivative is one.  This holds for every pair.

The family does not refute GU.  Indeed (3) and (6) give the exact restricted
class theorem
\[
 {G(q_n)\over I(q_n)}=3-2^{2-n}<3,qquad
 {U(q_n)\over I(q_n)}=1.                              \tag{12}
\]
Hence it satisfies GU with $(a,b)=(3,0)$ uniformly.  What fails is the
minimal-counterexample induction helper: restriction creates large negative
child $G$, so $\Delta_iG$ counts the same global reserve in every coordinate.
The next mechanism must control $G$ directly or telescope a nonreused signed
KL reserve; it cannot demand a bounded local residual derivative on all
tensor-prime Boolean functions.

